% Zone „Das kennst du schon“ – aus bank/prozentrechnung/zone.jsonl (zusammenbau v0.1)
\einheitenkopf[zone][Kennst du schon]{Das kennst du schon}
\setcounter{aufgabe}{0}

%% TODO Titel: Ich-kann-Satz fehlt in der Bank (Platzhalter: Kettenname) – Nr. 1
%% TODO Anweisung: ein Satz über der ersten Teilaufgabe fehlt in der Bank – Nr. 1
\begin{aufgabe}{Bruch als Anteil lesen und kürzen (drei von zwölf, gekürzt ein Viertel), auf den Nenner hundert erweitern}
\begin{teile}
\teil $6$ von $10$ Bonbons sind rot – welcher Anteil? Kürze. \leerfeld
\teil $6$ von $15$ Plätzen sind besetzt – welcher Anteil? Kürze und erweitere auf Hundertstel. \leerfeld
\end{teile}
\end{aufgabe}

%% TODO Titel: Ich-kann-Satz fehlt in der Bank (Platzhalter: Kettenname) – Nr. 2
%% TODO Anweisung: ein Satz über der ersten Teilaufgabe fehlt in der Bank – Nr. 2
\begin{aufgabe}{Bruch ↔ Dezimalzahl (ein Viertel, drei Fünftel)}
\begin{teile}
\teil $\frac{1}{2}$ als Dezimalzahl? \leerfeld
\teil $\frac{3}{25}$ als Dezimalzahl? \leerfeld
\end{teile}
\end{aufgabe}

%% TODO Titel: Ich-kann-Satz fehlt in der Bank (Platzhalter: Kettenname) – Nr. 3
%% TODO Anweisung: ein Satz über der ersten Teilaufgabe fehlt in der Bank – Nr. 3
\begin{aufgabe}{Durch hundert teilen, mit Dezimalzahlen multiplizieren (Kommaverschiebung)}
\begin{teile}
\teil $600 : 100$ \leerfeld
\teil $0{,}15 \cdot 60$ \leerfeld
\end{teile}
\end{aufgabe}

%% TODO Titel: Ich-kann-Satz fehlt in der Bank (Platzhalter: Kettenname) – Nr. 4
%% TODO Anweisung: ein Satz über der ersten Teilaufgabe fehlt in der Bank – Nr. 4
\begin{aufgabe}{Hoch- und Runterrechnen in einer Tabelle (Dreisatz, proportionale Zuordnung)}
\begin{teile}
\teil $2$ Kinokarten kosten $18\,€$ – wie viel kosten $5$ Karten?

\begin{dreisatz}{Karten}{Preis}\dsz{2}{18\,€}\dsp{\phantom{:0}}{\phantom{:0}}\dsz{1}{\dsleer}\dsp{\phantom{:0}}{\phantom{:0}}\dsz{5}{\dsleer}\end{dreisatz}
\teil $5$ Brötchen kosten $2\,€$ – wie viel kosten $3$ Brötchen?

\begin{dreisatz}{Brötchen}{Preis}\dsz{5}{2\,€}\dsp{\phantom{:0}}{\phantom{:0}}\dsz{1}{\dsleer}\dsp{\phantom{:0}}{\phantom{:0}}\dsz{3}{\dsleer}\end{dreisatz}
\end{teile}
\end{aufgabe}

%% TODO Titel: Ich-kann-Satz fehlt in der Bank (Platzhalter: Kettenname) – Nr. 5
%% TODO Anweisung: ein Satz über der ersten Teilaufgabe fehlt in der Bank – Nr. 5
\begin{aufgabe}{Runden auf eine Dezimalstelle}
\begin{teile}
\teil $4{,}73$ – auf eine Dezimalstelle? \leerfeld
\teil $18{,}456$ kg – auf eine Dezimalstelle? \leerfeld[kg]
\end{teile}
\end{aufgabe}

%% TODO Titel: Ich-kann-Satz fehlt in der Bank (Platzhalter: Kettenname) – Nr. 6
%% TODO Anweisung: ein Satz über der ersten Teilaufgabe fehlt in der Bank – Nr. 6
\begin{aufgabe}{Bruchteil einer Größe (zwei Drittel von 60 €)}
\begin{teile}
\teil $\frac{1}{4}$ von $20\,€$ \leerfeld[€]
\teil $\frac{3}{5}$ von $45$ m \leerfeld[m]
\end{teile}
\end{aufgabe}

%% TODO Titel: Ich-kann-Satz fehlt in der Bank (Platzhalter: Kettenname) – Nr. 7
%% TODO Anweisung: ein Satz über der ersten Teilaufgabe fehlt in der Bank – Nr. 7
\begin{aufgabe}{Bruchteil einer Größe (zwei Drittel von 60 €) – Fehler finden}
\begin{teile}
\teil Von $40\,€$ ist ein Fünftel ausgegeben – wie viel ist übrig? Ben rechnet: $40 : 5 = 8$, übrig sind $8\,€$. Wo ist der Fehler?
\end{teile}
\end{aufgabe}

%% TODO Titel: Ich-kann-Satz fehlt in der Bank (Platzhalter: Kettenname) – Nr. 8
%% TODO Anweisung: ein Satz über der ersten Teilaufgabe fehlt in der Bank – Nr. 8
\begin{aufgabe}{Bruchteil einer Größe (zwei Drittel von 60 €) – selbst rechnen}
\begin{teile}
\teil Von $30$ Litern im Tank ist ein Drittel verbraucht – wie viel ist übrig? \leerfeld[l]
\end{teile}
\end{aufgabe}
